\section*{Chapter \ref{ch:pl:native-extensions} --- Native Extensions}

\begin{solution}{exo:pl:native-extensions:1}
Per element through pybind11: $118 \times 2\times10^8$~ns $= 23.6$~s; in Python: 12.0~s; one call on the whole array at 1.4~ns: 0.28~s.
\end{solution}

\begin{solution}{exo:pl:native-extensions:2}
\texttt{ctypes} converts each argument at run time by the declared types, through generic machinery, and builds a C call frame
dynamically; pybind11's conversions are compiled into the binding for the exact signature. The machine instructions of the kernel are the
same; the work around them is not.
\end{solution}

\begin{solution}{exo:pl:native-extensions:3}
It converts the array into a new float64 array and passes the copy. For \texttt{ewma\_into} the kernel then writes into the copy, which is
discarded, and the caller's array is unchanged: a silent wrong result.
\end{solution}

\begin{solution}{exo:pl:native-extensions:4}
For $n > c_0/(p - c_1) = 118/58.6 \approx 2.0$: from three elements, one call on the whole array beats the Python loop.
\end{solution}

\begin{solution}{exo:pl:native-extensions:5}
numpy's lookup is a binary search for each left time, $O(\log m)$ steps each, and each step touches a different part of a larger array
(more cache misses as it grows); the merge walks both sorted arrays once, forward, a constant cost per element with sequential memory
access.
\end{solution}

\begin{solution}{exo:pl:native-extensions:6}
The array can be freed or reallocated by Python after the call (garbage-collected, resized, replaced), leaving the kernel with a dangling
pointer: reads of freed memory or writes into someone else's. The rule: a kernel reads and writes the buffers only during the call and keeps
no pointer afterwards; if data must persist, the caller keeps the array alive and passes it again.
\end{solution}

\begin{solution}{exo:pl:native-extensions:7}
All four agree on every one of the 100\,000 indices: $-1$ for the time before the first right time, and 49\,999 (the last index) for the
time after the last.
\end{solution}

\begin{solution}{exo:pl:native-extensions:8}
Converting a list of Python floats into a vector and a new list back is a copy and an object creation per element, the cost native code was
meant to avoid; a benchmark of ten numbers measures the call's fixed cost, not the per-element cost that matters on real data. Take and
return numpy arrays through the \omterm{def:pl:native-extensions:buffer}{buffer protocol}, and benchmark on realistic lengths.
\end{solution}

\begin{solution}{pb:pl:native-extensions:1}
\begin{enumerate}
\item It locates the function, checks and converts each argument from a Python object, may release and reacquire the lock, converts the
result back and checks for errors.
\item About 60~ns (a Python function), 118~ns (C++ through pybind11) and 337~ns (Rust through \texttt{ctypes}), loop included.
\item The Python loop around it remains, and each call adds conversions that cost more than the step itself.
\item Other Python threads run while the native code computes; it is safe only if the native code touches no Python object while the lock
is released.
\item The machine-level calling and layout convention; C's is simple, stable and spoken by every language's \omterm{def:pl:native-extensions:ffi}{foreign function interface}.
\item The argument's type, contiguity and dtype, and it converts or refuses accordingly.
\item That each array is a numpy array of the right dtype and C-contiguous, and pass its length.
\item Both refuse it: pybind11 because the argument is \texttt{c\_style} without \texttt{forcecast} and \texttt{noconvert}; the wrapper
because it checks contiguity.
\item It would copy silently: an output written into a copy is lost, and a copy of a large array costs what the extension was meant to
save.
\item Generated argument checks and conversions for Rust, like pybind11's for C++, and Rust types for Python objects; underneath, the same C
interface.
\item About 1.4~ns (C++), 1.5~ns (Rust) and 3.3~ns (numpy's filter).
\item The filter routine has a fixed cost of several microseconds per call, larger than the bindings', and a slower inner loop.
\item About 21~ns (numpy), 6~ns (C++) and 6~ns (Rust).
\item From about a thousand elements.
\item The lookup: its native version has a better algorithm (a merge instead of a binary search per element), not only a faster loop.
\item \textbf{Named result.} A call costs about 60~ns as a Python function, 118~ns through pybind11 and 337~ns through \texttt{ctypes}; on
whole arrays both native averages beat numpy at every length measured and the Rust lookup from about a thousand elements, and moving the
loop instead of the body turns a 24-second day of per-element calls into 0.3~s.
\item About 12~s against the Python loop (12.0~s to 0.28~s), 23~s against per-element native calls, and 0.4~s against numpy.
\item When the computation is already an array operation of a library, when arrays are short and calls are few, or when the gain does
not pay for building, testing and distributing compiled code on every platform the researchers use.
\item The kernel keeps no pointer after the call and frees nothing it did not allocate; the caller keeps its arrays alive and unchanged
during the call.
\item Around whole arrays: data cross once, loops stay native.
\end{enumerate}
\end{solution}

\begin{solution}{iq:pl:native-extensions:1}
It was called once per element from the Python loop: each call pays a conversion and dispatch cost larger than the arithmetic, and the
loop stayed in Python. Move the loop into the native function and pass whole arrays.
\iqlookfor{Per-call cost versus per-element cost.}
\end{solution}

\begin{solution}{iq:pl:native-extensions:2}
Through the \omterm{def:pl:native-extensions:buffer}{buffer protocol}: pass the array's pointer, dtype and shape (pybind11's \texttt{array\_t}, or \texttt{ctypes} pointers from
\texttt{arr.ctypes}). Wrong dtype or strides (read as if contiguous), silent copies (forcecast), writing to a read-only array, and pointers
kept after the call.
\iqlookfor{Layout checks and ownership.}
\end{solution}

\begin{solution}{iq:pl:native-extensions:3}
pybind11: compiled C++ bindings with type checks, conversions and lock handling, the norm for C++. \texttt{ctypes}: no compilation of a
binding, works with any C interface, but checks nothing and costs more per call. A Rust generator (PyO3): pybind11's comfort for Rust;
the C interface with \texttt{ctypes} works without it.
\iqlookfor{Safety and per-call cost against convenience and build complexity.}
\end{solution}

\begin{solution}{iq:pl:native-extensions:4}
So that other Python threads run while it computes. It must not touch any Python object, including reference counts, raise Python
exceptions, or call back into Python until it has reacquired the lock.
\iqlookfor{No Python API without the lock.}
\end{solution}

\begin{solution}{iq:pl:native-extensions:5}
Kernels chosen by profiling (loops that arrays cannot express, merges, state machines); written in C++ or Rust against plain pointers, tested
alone; bound once per array with pybind11 (or a C interface), refusing layouts they cannot read in place and releasing the lock; built in
the continuous-integration pipeline for each supported platform and Python version and distributed as wheels; tested against a Python
reference on the same data.
\iqlookfor{Selection by measurement, separation of kernel and binding, reproducible builds.}
\end{solution}
